\chapter{Rings and Ideals}

\begin{center}\itshape
In $\Z[\sqrt{-5}]$, $6=2\cdot3=(1+\sqrt{-5})(1-\sqrt{-5})$, and none of these four numbers can be split further.\\
Unique factorisation, which we took for granted in $\Z$, simply fails here. What property of $\Z$ was doing the work?
\end{center}

In the integers, $6=2\cdot3$ is a prime factorisation, and every other factorisation agrees with it up to signs and order. But in the larger ring
\[
\Z[\sqrt{-5}]=\{a+b\sqrt{-5}:a,b\in\Z\}
\]
we also have
\[
6=(1+\sqrt{-5})(1-\sqrt{-5}).
\]
These two expressions are not the same factorisation in disguise. The usual uniqueness of prime factors has failed.

That failure is a reason to study \emph{rings}. A ring has addition, subtraction and multiplication, but division may be impossible. Polynomials form rings. So do residue classes modulo $n$, the Gaussian integers, and many sets of algebraic numbers. Their common rules give us tools that work in all these examples.

An ideal is to a ring what a normal subgroup is to a group: the right kind of subset to collapse when building a quotient. There is an extra demand. Multiplying an element of an ideal by \emph{any} element of the ring must keep it inside the ideal. That makes multiplication well defined in the quotient.

Some ideals behave like prime numbers; others give a field when we divide by them. Those two properties become the tests for prime and maximal ideals. Finally, Gauss's lemma tells us when familiar polynomial factorisation over a field can be brought back to the original ring.

The failed factorisation of $6$ will return in arithmetic. Here it teaches us what algebra can still say when ordinary prime factorisation is no longer guaranteed.

\section{Rings}

\begin{definition}
A (commutative) \emph{ring} is a set $R$ with two operations $+$ and $\times$ such that $(R,+)$ is an abelian group with neutral element $0$, multiplication is associative and commutative with a neutral element $1$, and $a(b+c)=ab+ac$. A \emph{ring homomorphism} $f:R\to S$ satisfies $f(a+b)=f(a)+f(b)$, $f(ab)=f(a)f(b)$ and $f(1)=1$. The \emph{units} of $R$ are the elements with a multiplicative inverse; they form a group $R^*$.

$R$ is an \emph{integral domain} (or \emph{domain}) if $1\ne0$ and $ab=0$ implies $a=0$ or $b=0$. It is a \emph{field} if $1\ne0$ and $R^*=R\setminus\{0\}$.
\end{definition}

All rings in this chapter are commutative. Non-commutative rings are important too ($n\times n$ matrices, and the quaternions of Chapter~12), but factorisation is a commutative story.

\begin{example}\label{ex:rings}
$\Z$, $\Q$, $\R$, $\C$; $\Z/n\Z$, a field iff $n$ is prime (Volume~2); $K[X]$ for a field $K$, and $R[X]$ for any ring $R$; $\Z[i]$, $\Z[\sqrt{-5}]=\{a+b\sqrt{-5}:a,b\in\Z\}$; the product $R\times S$ with componentwise operations, never a domain (as $(1,0)(0,1)=0$); the continuous functions $\R\to\R$. A field is a domain: if $ab=0$ and $a\ne0$, then $b=a^{-1}ab=0$. A domain may be far from a field: $\Z$, $K[X]$.
\end{example}

\begin{proposition}\label{prop:finitedomain}
A finite integral domain is a field.
\end{proposition}
\begin{proof}
Let $a\ne0$. The map $x\mapsto ax$ is injective, since $ax=ay$ gives $a(x-y)=0$, so $x=y$. An injective map from a finite set to itself is onto, so $ax=1$ for some $x$.
\end{proof}

In a domain one may cancel: $ab=ac$ with $a\ne0$ gives $b=c$. If $R$ is a domain, so is $R[X]$: the leading coefficient of $fg$ is the product of the leading coefficients, so $\deg(fg)=\deg f+\deg g$ for $f,g\ne0$. In particular $R[X]^*=R^*$.

A ring homomorphism may send several elements to the same output. As with groups, the kernel measures what the map forgets. But multiplication puts a new demand on the kernel: it must absorb products by arbitrary ring elements.

\section{Ideals and quotient rings}

\begin{definition}
An \emph{ideal} of $R$ is an additive subgroup $I$ such that $ra\in I$ for all $r\in R$, $a\in I$. The ideal \emph{generated} by $a_1,\dots,a_n$ is $(a_1,\dots,a_n)=\{r_1a_1+\dots+r_na_n:r_i\in R\}$; an ideal $(a)=aR$ is \emph{principal}.
\end{definition}

The kernel of a ring homomorphism $f$ is an ideal: if $f(a)=0$ then $f(ra)=f(r)f(a)=0$. An ideal containing a unit $u$ contains $u^{-1}u=1$ and then everything: $I=R$. So a field $K$ has only the ideals $\{0\}$ and $K$, and conversely a ring $R$ with $1\ne0$ whose only ideals are $\{0\}$ and $R$ is a field, since for $a\ne0$ the ideal $(a)$ must be $R$, so $ab=1$ for some $b$.

\begin{proposition}[quotient rings]\label{prop:quotring}
Let $I\ne R$ be an ideal.
\begin{enumerate}[label=(\alph*)]
\item The quotient group $R/I$ is a ring with $(a+I)(b+I)=ab+I$, and $\pi:a\mapsto a+I$ is a ring homomorphism with kernel $I$.
\item (First isomorphism theorem) If $f:R\to S$ is a ring homomorphism, $R/\ker f\cong f(R)$.
\item (Correspondence) $J\mapsto J/I$ is a bijection from the ideals of $R$ containing $I$ to the ideals of $R/I$.
\end{enumerate}
\end{proposition}
\begin{proof}
(a) $(R,+)$ is abelian, so $I$ is a normal subgroup and $R/I$ is an abelian group (Theorem~\ref{thm:quotient}). The product is well defined: if $a'=a+i$, $b'=b+j$ with $i,j\in I$, then $a'b'=ab+(aj+ib+ij)\in ab+I$. The axioms hold in $R/I$ because they hold in $R$, and $\pi$ respects all operations by construction.

(b) The group isomorphism $a+\ker f\mapsto f(a)$ of Theorem~\ref{thm:iso1} also respects products and $1$.

(c) By Theorem~\ref{thm:iso23}(a), $J\mapsto\pi(J)$ is a bijection between the additive subgroups containing $I$ and those of $R/I$, with inverse $\pi^{-1}$. If $J$ is an ideal, $(r+I)(a+I)=ra+I\in\pi(J)$, so $\pi(J)$ is an ideal; if $J'$ is an ideal of $R/I$, $\pi(ra)=\pi(r)\pi(a)\in J'$ for $a\in\pi^{-1}(J')$, so $\pi^{-1}(J')$ is an ideal.
\end{proof}

\begin{example}\label{ex:quotients}
\begin{enumerate}[label=(\arabic*)]
\item The ideals of $\Z$ are the $n\Z$ (they are subgroups), and $\Z/n\Z$ is the ring of Volume~2.
\item Evaluation at $a$, $K[X]\to K$, $P\mapsto P(a)$, is onto with kernel $(X-a)$ (a polynomial vanishes at $a$ iff $X-a$ divides it): $K[X]/(X-a)\cong K$.
\item Evaluation at $i$, $\R[X]\to\C$, is onto ($a+bX\mapsto a+bi$) with kernel $(X^2+1)$: a real polynomial vanishing at $i$ vanishes at $-i$, so it is divisible by $(X-i)(X+i)$. Hence $\R[X]/(X^2+1)\cong\C$. In the same way $\Z[X]/(X^2+1)\cong\Z[i]$ and $\Z[X]/(X^2+5)\cong\Z[\sqrt{-5}]$ (division by a monic polynomial stays in $\Z[X]$).
\end{enumerate}
\end{example}

So $\C$ can be \emph{defined} as $\R[X]/(X^2+1)$: adjoin a symbol $X$ and decree $X^2+1=0$. This is how all field extensions will be built in Chapter~8.

The quotient by an ideal is always a ring. Sometimes it has no zero divisors; sometimes every non-zero element has an inverse. Those two outcomes tell us exactly which ideals deserve the names prime and maximal.

\section{Prime and maximal ideals}

\begin{definition}
An ideal $P\ne R$ is \emph{prime} if $ab\in P$ implies $a\in P$ or $b\in P$. An ideal $M\ne R$ is \emph{maximal} if the only ideals containing $M$ are $M$ and $R$.
\end{definition}

\begin{proposition}\label{prop:primemax}
Let $I\ne R$ be an ideal.
\begin{enumerate}[label=(\alph*)]
\item $I$ is prime iff $R/I$ is an integral domain.
\item $I$ is maximal iff $R/I$ is a field.
\item Every maximal ideal is prime.
\end{enumerate}
\end{proposition}
\begin{proof}
(a) $(a+I)(b+I)=0$ in $R/I$ means $ab\in I$, and $a+I=0$ means $a\in I$: the two definitions say the same thing. (And $1\ne0$ in $R/I$ because $I\ne R$.)

(b) $R/I$ is a field iff its only ideals are $0$ and $R/I$, iff (Proposition~\ref{prop:quotring}(c)) the only ideals containing $I$ are $I$ and $R$.

(c) A field is a domain.
\end{proof}

\paragraph{Why this works.}
There are two different questions about a quotient ring. Can two non-zero classes multiply to zero? That happens precisely when $ab\in I$ but neither $a$ nor $b$ is in $I$, which is what a prime ideal forbids. Can we divide by every non-zero class? That happens precisely when the quotient is a field, or equivalently when it has no non-trivial ideals. For a concrete contrast, $(6)$ is not prime in $\Z$, since $2\cdot3\in(6)$ but neither factor lies in $(6)$. The ideal $(5)$ is maximal: $\Z/(5)$ is a field. In $\Z$, a non-zero prime ideal is maximal; in $\Z[X]$ this need not hold.

In $\Z$, the prime ideals are $\{0\}$ and the $p\Z$, and the non-zero ones are maximal. In $\Z[X]$ the picture is richer:

\begin{think}
Is $(X)$ a maximal ideal of $\Z[X]$? Is $(2,X)$ principal? Compute $\Z[X]/(X)$ and $\Z[X]/(2,X)$ before reading on.
\end{think}

$\Z[X]/(X)\cong\Z$ (evaluate at $0$) is a domain but not a field: $(X)$ is prime, not maximal. $\Z[X]/(2,X)\cong\Z/2\Z$ (evaluate at $0$, then reduce mod $2$): $(2,X)$ is maximal. And $(2,X)$ is not principal (Exercise~\ref{ex:zxnotpid}).

\begin{theorem}[Chinese remainder theorem for rings]\label{thm:crtring}
Let $I,J$ be ideals with $I+J=R$ (they are \emph{coprime}). Then $I\cap J=IJ$, where $IJ$ is the ideal generated by the products $ab$, $a\in I$, $b\in J$, and
\[
\begin{gathered}R/(I\cap J)\cong R/I\times R/J,\qq x+I\cap J\mapsto(x+I,\,x+J).\end{gathered}
\]
If $I_1,\dots,I_n$ are pairwise coprime, $R/(I_1\cap\dots\cap I_n)\cong\prod_kR/I_k$.
\end{theorem}
\begin{proof}
Write $1=a+b$ with $a\in I$, $b\in J$. Always $IJ\subseteq I\cap J$; conversely, for $x\in I\cap J$, $x=xa+xb\in IJ$. The homomorphism $\varphi:R\to R/I\times R/J$, $x\mapsto(x+I,x+J)$, has kernel $I\cap J$. It is onto: $b=1-a\equiv1\pmod I$ and $b\equiv0\pmod J$, and $a\equiv0\pmod I$, $a\equiv1\pmod J$, so $x=ub+va$ maps to $(u+I,v+J)$. Conclude with Proposition~\ref{prop:quotring}(b).

For $n$ ideals, by induction it suffices that $I_1$ and $I_2\cap\dots\cap I_n$ are coprime. Write $1=a_k+b_k$ with $a_k\in I_1$, $b_k\in I_k$ for $k\ge2$; then $1=\prod_{k\ge2}(a_k+b_k)$, and expanding, every term except $b_2\cdots b_n\in I_2\cap\dots\cap I_n$ contains some $a_k$, so lies in $I_1$.
\end{proof}

For $R=\Z$ this is the theorem of Volume~2. For $R=K[X]$ and $I_k=(X-x_k)$ with distinct $x_k$, it says that a polynomial is determined modulo $\prod(X-x_k)$ by its values at the $x_k$, and that any values can be prescribed: this is Lagrange interpolation (Volume~1, Chapter~7), and the element $b$ of the proof is a Lagrange basis polynomial.

A quotient field is the strongest possible setting for division. Inside a domain, however, we may not be able to divide, and factorisation can behave very differently. We need to separate the properties that integers happen to enjoy from those that hold in every domain.

\section{Factorisation in domains}

From now on $R$ is an integral domain. We write $a\mid b$ if $b=ac$ for some $c\in R$; this means $(b)\subseteq(a)$. Elements $a,b$ are \emph{associates} if $a=ub$ with $u\in R^*$, equivalently $(a)=(b)$.

\begin{definition}
A non-zero non-unit $p$ is \emph{irreducible} if $p=ab$ implies that $a$ or $b$ is a unit. It is \emph{prime} if $p\mid ab$ implies $p\mid a$ or $p\mid b$, i.e.\ if $(p)$ is a prime ideal. $R$ is a \emph{unique factorisation domain} (UFD) if every non-zero non-unit is a product of irreducibles, uniquely up to order and associates.
\end{definition}

A prime $p$ is irreducible: if $p=ab$, then $p\mid a$ say, $a=pc$, so $p=pcb$ and $cb=1$: $b$ is a unit. In $\Z$, ``prime'' and ``irreducible'' coincide, which is Euclid's lemma. The opening example shows how they can differ:

\begin{proposition}\label{prop:zm5}
In $\Z[\sqrt{-5}]$, with the norm $N(a+b\sqrt{-5})=a^2+5b^2=|a+b\sqrt{-5}|^2$:
\begin{enumerate}[label=(\alph*)]
\item $N(xy)=N(x)N(y)$, the units are $\pm1$, and no element has norm $2$ or $3$;
\item $2$, $3$, $1+\sqrt{-5}$ and $1-\sqrt{-5}$ are irreducible, and no two of them are associates;
\item $2$ is irreducible but not prime. So $\Z[\sqrt{-5}]$ is not a UFD.
\end{enumerate}
\end{proposition}
\begin{proof}
(a) The norm is the square of the complex modulus, which is multiplicative. If $xy=1$ then $N(x)N(y)=1$, so $N(x)=1$, i.e.\ $a^2+5b^2=1$, so $b=0$, $a=\pm1$. And $a^2+5b^2\in\{2,3\}$ forces $b=0$ and $a^2\in\{2,3\}$, impossible (Figure~\ref{fig:lattice}).

(b) The norms are $4,9,6,6$. If one of these elements were $xy$ with $x,y$ non-units, $N(x)$ would be a proper divisor of $4$, $9$ or $6$ different from $1$, so $2$ or $3$: impossible by (a). Associates have the same norm, and $1+\sqrt{-5}\ne\pm(1-\sqrt{-5})$.

(c) $2$ divides $6=(1+\sqrt{-5})(1-\sqrt{-5})$ but divides neither factor, since $\frac{1\pm\sqrt{-5}}2\notin\Z[\sqrt{-5}]$. So $6$ has two factorisations into irreducibles that differ by more than order and units.
\end{proof}

\begin{figure}[ht]
\centering
\fig{ch06-lattice}
\caption{$\Z[\sqrt{-5}]$ as the lattice of points $a+b\sqrt5\,i$. The norm is the squared distance to $0$. The circles of norm $2$ and $3$ (dashed) contain no lattice point, so nothing can divide $2$, $3$ or $1\pm\sqrt{-5}$ properly, and $6$ factors in two different ways. In Volume~13, Chapter~2, the ideal $(2,1+\sqrt{-5})$ restores uniqueness: $(2)=(2,1+\sqrt{-5})^2$.}
\label{fig:lattice}
\end{figure}

\begin{lemma}\label{lem:primeufd}
In a domain, if $a=p_1\cdots p_r=q_1\cdots q_s$ with the $p_i$ prime and the $q_j$ irreducible, then $r=s$ and, after reordering, $q_i$ is an associate of $p_i$. Hence a domain in which every non-zero non-unit is a product of irreducibles is a UFD iff every irreducible is prime.
\end{lemma}
\begin{proof}
Induction on $r$. If $r=0$, $a=1$ and the $q_j$, being non-units, cannot be there: $s=0$. If $r\ge1$, $p_1$ divides $q_1\cdots q_s$, so it divides some $q_j$, say $q_1=p_1c$; as $q_1$ is irreducible and $p_1$ is not a unit, $c$ is a unit. Cancel $p_1$: $p_2\cdots p_r=(cq_2)q_3\cdots q_s$, and apply the induction hypothesis.

For the last sentence: if every irreducible is prime, the first part gives uniqueness. Conversely, in a UFD let $p$ be irreducible with $p\mid ab$, $ab=pc$. Factor $a$, $b$, $c$ into irreducibles; by uniqueness, the irreducible $p$ is an associate of one of the factors of $a$ or of $b$, so $p\mid a$ or $p\mid b$. (If $a$ or $b$ is $0$ or a unit, this is immediate.)
\end{proof}

\begin{definition}
$R$ is a \emph{principal ideal domain} (PID) if every ideal is principal. $R$ is \emph{Euclidean} if there is a map $N:R\setminus\{0\}\to\N$ such that for all $a$ and all $b\ne0$ there are $q,r$ with $a=bq+r$ and $r=0$ or $N(r)<N(b)$.
\end{definition}

$\Z$ ($N=|\cdot|$), $K[X]$ ($N=\deg$), $\Z[i]$ and $\Z[\omega]$ ($N=|\cdot|^2$, Volume~2, Chapter~10) are Euclidean.

\begin{theorem}\label{thm:edpidufd}
Euclidean $\Rightarrow$ PID $\Rightarrow$ UFD.
\end{theorem}
\begin{proof}
\emph{Euclidean $\Rightarrow$ PID.} Let $I\ne\{0\}$ be an ideal and $b\in I\setminus\{0\}$ with $N(b)$ minimal. For $a\in I$, write $a=bq+r$; then $r=a-bq\in I$, and $N(r)<N(b)$ is impossible, so $r=0$ and $a\in(b)$. Thus $I=(b)$.

\emph{PID $\Rightarrow$ UFD.} First, every irreducible $p$ is prime. If $(p)\subseteq(d)$ then $p=dc$, so $d$ is a unit ($(d)=R$) or $c$ is a unit ($(d)=(p)$). So $(p)$ is maximal (every ideal is some $(d)$), hence prime (Proposition~\ref{prop:primemax}).

Next, every non-zero non-unit $a$ is a product of irreducibles. Suppose not. Then $a$ is not irreducible, so $a=a_1b_1$ with non-units $a_1,b_1$, and one of them, say $a_1$, is not a product of irreducibles either. Repeating gives a sequence with $(a)\subsetneq(a_1)\subsetneq(a_2)\subsetneq\cdots$ (strict, since $b_1$ is not a unit). The union of an increasing chain of ideals is an ideal, so it is $(d)$ for some $d$; then $d$ belongs to some $(a_k)$, and $(a_k)\subseteq(a_{k+1})\subseteq(d)\subseteq(a_k)$, contradicting strictness. Conclude with Lemma~\ref{lem:primeufd}.
\end{proof}

In a PID, $\gcd$s exist and Bézout holds: $(a,b)=(d)$ means that $d$ divides $a$ and $b$ and $d=ua+vb$. The converse implications fail: $\Z[\frac{1+\sqrt{-19}}2]$ is a PID that is not Euclidean, and $\Z[X]$ is a UFD (Theorem~\ref{thm:gaussufd}) that is not a PID (Exercise~\ref{ex:zxnotpid}).

The integers have unique factorisation even though division is not always possible. We now ask whether that property survives when we adjoin a variable. The key is to distinguish the factors that come from coefficients from those carried by the polynomial itself.

\section{Polynomials over a UFD}

For a domain $R$, the \emph{field of fractions} $K=\operatorname{Frac}R$ is built exactly as $\Q$ is built from $\Z$: fractions $a/b$ with $b\ne0$, where $a/b=c/d$ iff $ad=bc$, with the usual operations; the verifications use only that $R$ is a domain, and $R\subseteq K$ via $a\mapsto a/1$. For example $\operatorname{Frac}\Z[i]=\Q(i)$ and $\operatorname{Frac}K[X]=K(X)$, the rational functions.

Let $R$ be a UFD. The \emph{content} $c(f)$ of a non-zero $f\in R[X]$ is a gcd of its coefficients (defined up to a unit), and $f$ is \emph{primitive} if $c(f)$ is a unit. Every non-zero $f\in K[X]$ can be written $f=\lambda f_0$ with $\lambda\in K^*$ and $f_0\in R[X]$ primitive: multiply by a common denominator, then divide by the content. This writing is unique up to a unit of $R$: if $\lambda f_0=\mu g_0$, write $\mu/\lambda=a/b$ in lowest terms (possible in a UFD); then $bf_0=ag_0$, comparing contents $b\sim a$ up to units, so $a/b\in R^*$.

\begin{lemma}[Gauss]\label{lem:gauss}
Let $R$ be a UFD. A product of primitive polynomials is primitive; hence $c(fg)=c(f)c(g)$ up to units.
\end{lemma}
\begin{proof}
First suppose that $f,g\in R[X]$ are primitive. We want to show that $fg$ is primitive as well.

Assume the contrary. Then the coefficients of $fg$ have a common non-unit divisor. Since $R$ is a UFD, that divisor has an irreducible factor $p$. Thus $p$ divides \emph{every} coefficient of $fg$. Reduce each coefficient modulo the ideal $(p)$. In $(R/(p))[X]$ this says
\[
\bar f\,\bar g=\overline{fg}=0.
\]
Neither $\bar f$ nor $\bar g$ is zero: otherwise $p$ would divide every coefficient of $f$ or of $g$, contradicting primitivity.

In a UFD every irreducible element is prime (Lemma~\ref{lem:primeufd}). Hence $(p)$ is a prime ideal, $R/(p)$ is an integral domain, and so is its polynomial ring. Two non-zero polynomials over a domain cannot have zero product. This contradiction proves that $fg$ is primitive.

For arbitrary non-zero polynomials write $f=c(f)f_0$ and $g=c(g)g_0$, where $f_0$ and $g_0$ are primitive. Then
\[
fg=c(f)c(g)\,(f_0g_0).
\]
We just proved that $f_0g_0$ is primitive, so the content of $fg$ is associated to $c(f)c(g)$. This is the stated identity up to multiplication by a unit of $R$.
\end{proof}

\paragraph{Why this works.}
The word \emph{primitive} means that no non-unit of the coefficient ring divides every coefficient. To test whether a product has gained a common prime factor $p$, reduce the coefficients modulo $p$. If $p$ divides every coefficient of $fg$, the product becomes the zero polynomial. But the reductions of primitive $f$ and $g$ are both non-zero, and a polynomial ring over a domain has no zero divisors. This contradiction explains the lemma. The UFD hypothesis matters: it lets us choose an irreducible $p$ dividing a non-unit common factor, and it makes that $p$ prime.

\begin{proposition}\label{prop:gaussirr}
Let $R$ be a UFD with fraction field $K$, and $f\in R[X]$ primitive of degree $\ge1$. Then $f$ is irreducible in $R[X]$ iff it is irreducible in $K[X]$. If $f$ divides $g\in R[X]$ in $K[X]$, it divides $g$ in $R[X]$.
\end{proposition}
\begin{proof}
If $f=GH$ in $K[X]$ with $\deg G,\deg H\ge1$, write $G=\lambda g$, $H=\mu h$ with $g,h$ primitive. Then $f=\lambda\mu\,gh$ with $gh$ primitive (Lemma~\ref{lem:gauss}), so by the uniqueness above $\lambda\mu\in R^*$, and $f=(\lambda\mu g)h$ is a factorisation in $R[X]$ into polynomials of degree $\ge1$, which are not units. Conversely, if $f=gh$ in $R[X]$ with $g,h$ non-units, neither can be a constant (a constant factor divides $c(f)$, a unit), so both have degree $\ge1$ and $f$ is reducible in $K[X]$.

If $g=fk$ with $k\in K[X]$, write $k=\lambda k_0$ with $k_0$ primitive. Then $g=\lambda\,fk_0$ with $fk_0$ primitive, so $\lambda$ is, up to a unit, the content of $g$, which lies in $R$. So $k\in R[X]$.
\end{proof}

\begin{theorem}[Gauss]\label{thm:gaussufd}
If $R$ is a UFD, so is $R[X]$. In particular $\Z[X]$ and $K[X_1,\dots,X_n]$ are UFDs.
\end{theorem}
\begin{proof}
\emph{Existence.} Let $f\in R[X]$ be a non-zero non-unit; write $f=c(f)f_0$. Factor $c(f)$ into irreducibles of $R$ (they stay irreducible in $R[X]$, by degree). Factor $f_0$ into irreducibles of the PID $K[X]$, $f_0=G_1\cdots G_k$, and write $G_i=\lambda_ig_i$ with $g_i$ primitive. As in the proof of Proposition~\ref{prop:gaussirr}, $f_0=u\,g_1\cdots g_k$ with $u\in R^*$, and each $g_i$ is irreducible in $R[X]$ by that proposition.

\emph{Irreducibles are prime.} An irreducible of $R[X]$ is either a constant $p$, irreducible in $R$, or a primitive polynomial irreducible in $K[X]$. In the first case, $R[X]/(p)\cong(R/(p))[X]$ (reduce the coefficients) is a domain, so $p$ is prime. In the second, if $f\mid gh$ in $R[X]$, then in the PID $K[X]$, where $f$ is irreducible hence prime, $f\mid g$ or $f\mid h$; by Proposition~\ref{prop:gaussirr} this holds in $R[X]$. Conclude with Lemma~\ref{lem:primeufd}. The case of $n$ variables follows by induction, as $K[X_1,\dots,X_n]=K[X_1,\dots,X_{n-1}][X_n]$.
\end{proof}

Volume~1, Chapter~7, proved Gauss' lemma and Eisenstein's criterion for $\Z[X]$; they now hold over any UFD, with the same proofs. Another tool is reduction: if $f\in\Z[X]$ is monic and its reduction mod $p$ is irreducible over $\Z/p\Z$, then $f$ is irreducible over $\Q$ (a factorisation into monic polynomials of lower degree would reduce to one mod $p$). The converse fails badly, see Exercise~\ref{ex:x4p1}.

\begin{beauty}
Kummer's discovery that factorisation fails in rings like $\Z[\zeta_{23}]$ broke Lamé's 1847 ``proof'' of Fermat's Last Theorem, and the repair, ideals, is the subject of Volume~13, Chapters~1 and~2: in the ring of integers of a number field every non-zero ideal factors uniquely into prime ideals, even when numbers do not. In computer algebra, the UFD property of $\Z[X]$ and $\Z[X,Y]$ is what makes polynomial gcd and factorisation algorithms possible, and the Chinese remainder theorem for rings is the basis of \emph{multi-modular} algorithms.
\end{beauty}

\section{Exercises}

\begin{exercise}[$\star$]\label{ex:idealsZn}
Describe the ideals of $\Z/n\Z$, and say which are prime and which are maximal. Find the units and the zero divisors of $\Z/12\Z$.
\end{exercise}

\begin{exercise}[$\star$]\label{ex:quot}
Show that $\R[X]/(X^2-1)\cong\R\times\R$, and that $\R[X]/(X^2)$ is not isomorphic to it, nor to $\C$. (Look for elements with $x^2=0$.)
\end{exercise}

\begin{exercise}[$\star\star$]\label{ex:zsqrtm5}
In $\Z[\sqrt{-5}]$, show that the ideal $I=(2,1+\sqrt{-5})$ is not principal, that $\Z[\sqrt{-5}]/I\cong\Z/2\Z$, so $I$ is maximal, and that $3$ is irreducible but not prime.
\end{exercise}

\begin{exercise}[$\star\star$]\label{ex:zi}
\begin{enumerate}[label=(\alph*)]
\item Show that every complex number $z$ is at distance at most $\frac1{\sqrt2}$ from some Gaussian integer, and deduce again that $\Z[i]$ is Euclidean for $N(z)=|z|^2$.
\item Do the same for $\Z[\sqrt{-2}]$.
\item Why does the argument fail for $\Z[\sqrt{-5}]$? Find $a,b$ for which no $q\in\Z[\sqrt{-5}]$ satisfies $N(a-bq)<N(b)$.
\end{enumerate}
\end{exercise}

\begin{exercise}[$\star\star$]\label{ex:zxnotpid}
Show that $(2,X)$ is not a principal ideal of $\Z[X]$, so $\Z[X]$ is not a PID. More generally show that $R[X]$ is a PID iff $R$ is a field.
\end{exercise}

\begin{exercise}[$\star\star$]\label{ex:boolean}
A ring is \emph{Boolean} if $x^2=x$ for every $x$. Show that in a Boolean ring $2x=0$ for all $x$, that the only Boolean domain is $\Z/2\Z$, and that every prime ideal of a Boolean ring is maximal. (Example: the subsets of a set, with symmetric difference as sum and intersection as product.)
\end{exercise}

\begin{exercise}[$\star\star$]\label{ex:crtinterp}
Use Theorem~\ref{thm:crtring} in $\Q[X]$ to find the polynomial $P$ of degree at most $2$ with $P(0)=1$, $P(1)=0$ and $P'(1)=2$. (Note that $P(1)=a$, $P'(1)=b$ is a condition modulo $(X-1)^2$.)
\end{exercise}

\begin{exercise}[$\star\star\star$]\label{ex:x4p1}
Show that $X^4+1$ is irreducible in $\Q[X]$, but that its reduction modulo $p$ is reducible for every prime $p$.
\end{exercise}

\section{Hints}

\begin{hint}{ex:idealsZn}
Use Proposition~\ref{prop:quotring}(c) with $R=\Z$ and $I=n\Z$.
\end{hint}
\begin{hint}{ex:quot}
Apply Theorem~\ref{thm:crtring} to $(X-1)$ and $(X+1)$.
\end{hint}
\begin{hint}{ex:zsqrtm5}
If $I=(d)$, then $N(d)$ divides $N(2)=4$ and $N(1+\sqrt{-5})=6$. For the quotient, $a+b\sqrt{-5}\equiv a+b\pmod I$. For $3$, look at $(1+\sqrt{-5})(1-\sqrt{-5})$ again.
\end{hint}
\begin{hint}{ex:zi}
(a) Round the real and imaginary parts of $a/b$. (c) Compute the farthest point of the plane from the lattice $\Z[\sqrt{-5}]$.
\end{hint}
\begin{hint}{ex:zxnotpid}
A generator of $(2,X)$ would divide $2$ and $X$. If $R$ is not a field, take a non-zero non-unit $a$ and look at $(a,X)$.
\end{hint}
\begin{hint}{ex:boolean}
Expand $(x+x)^2$ and $(x+y)^2$. In a domain, $x(x-1)=0$. For prime ideals, look at $R/P$.
\end{hint}
\begin{hint}{ex:crtinterp}
$(X)$ and $(X-1)^2$ are coprime. Find polynomials congruent to $1$ modulo one and to $0$ modulo the other.
\end{hint}
\begin{hint}{ex:x4p1}
Over $\Q$, substitute $X\mapsto X+1$. Modulo an odd prime $p$, one of $-1$, $2$, $-2$ is a square (why?); write $X^4+1$ as a difference of two squares in three different ways.
\end{hint}

\section{Solutions}

\begin{solution}{ex:idealsZn}
\idea Ideals of a quotient are ideals of $\Z$ containing $n\Z$, and those are the $d\Z$ with $d\mid n$.
\writeup By Proposition~\ref{prop:quotring}(c), the ideals of $\Z/n\Z$ are the $d\Z/n\Z$ with $n\Z\subseteq d\Z$, i.e.\ $d\mid n$ ($d>0$). The quotient is $(\Z/n\Z)/(d\Z/n\Z)\cong\Z/d\Z$ (third isomorphism theorem), which is a domain iff it is a field iff $d$ is prime. So the prime ideals and the maximal ideals are the same: the $p\Z/n\Z$ for $p\mid n$ prime.

In $\Z/12\Z$, $a$ is a unit iff $\gcd(a,12)=1$: $1,5,7,11$. Every other non-zero element is a zero divisor: if $d=\gcd(a,12)>1$ then $a\cdot\frac{12}d\equiv0$ with $\frac{12}d\not\equiv0$. So $2,3,4,6,8,9,10$ are zero divisors. In a finite ring every element is a unit, a zero divisor, or $0$.
\end{solution}

\begin{solution}{ex:quot}
\idea Coprime factors split the quotient into a product (CRT); a repeated factor creates a non-zero element whose square is $0$.
\writeup $(X-1)+(X+1)$ contains $\frac12\bigl((X+1)-(X-1)\bigr)=1$, so the two ideals are coprime, with product $(X^2-1)$. By Theorem~\ref{thm:crtring}, $\R[X]/(X^2-1)\cong\R[X]/(X-1)\times\R[X]/(X+1)\cong\R\times\R$ (Example~\ref{ex:quotients}). In $\R[X]/(X^2)$ the class $x$ of $X$ is non-zero with $x^2=0$. In $\R\times\R$, $(a,b)^2=0$ forces $a=b=0$, and in the field $\C$ too; an isomorphism would carry $x$ to such an element. So $\R[X]/(X^2)$ is isomorphic to neither.
\end{solution}

\begin{solution}{ex:zsqrtm5}
\idea Three tools handle the three requests. Use the integer norm to test whether the ideal could have one generator; reduce coefficients modulo $2$ to recognise the quotient; use the factorisation of $6$ to show that an irreducible element need not be prime.

\writeup Put $R=\Z[\sqrt{-5}]$ and $I=(2,1+\sqrt{-5})$. Every element of $R$ is uniquely $a+b\sqrt{-5}$ with $a,b\in\Z$. Its norm
\[
N(a+b\sqrt{-5})=a^2+5b^2
\]
is a non-negative integer and is multiplicative.

\emph{The ideal is not principal.} Suppose $I=(d)$. Then $d$ divides both $2$ and $1+\sqrt{-5}$, so the positive integer $N(d)$ divides
\[
N(2)=4,\qquad N(1+\sqrt{-5})=6.
\]
Thus $N(d)=1$ or $2$. There is no element of norm $2$: the equation $a^2+5b^2=2$ has no integer solution. Hence $N(d)=1$ and $d$ is a unit, forcing $I=R$.

But $I$ is proper. Indeed, the map
\[
\psi:R\to\mathbb F_2,\qquad
\psi(a+b\sqrt{-5})=[a+b]
\]
respects multiplication: modulo $2$, the relation $(\sqrt{-5})^2=-5$ becomes $1$, so sending $\sqrt{-5}$ to $1$ is compatible with multiplication. Both generators of $I$ map to zero, while $\psi(1)=1$. Consequently $I\subsetneq R$, a contradiction. So $I$ is not principal.

\emph{The quotient is a field.} The map $\psi$ is onto because $\psi(1)=1$. Its kernel consists of elements with $a+b$ even. Such an element can be written
\[
a+b\sqrt{-5}=(a-b)+b(1+\sqrt{-5}),
\]
and $a-b$ is even precisely when $a+b$ is. Thus $\ker\psi=I$. The first isomorphism theorem for rings gives
\[
R/I\cong\mathbb F_2.
\]
Since this quotient is a field, $I$ is maximal.

\emph{Irreducible is not prime.} An element of norm $3$ does not exist: $a^2+5b^2=3$ has no integer solution. If $3=uv$ with both factors non-units, then $N(u)N(v)=N(3)=9$ and both norms are integers greater than $1$. The only possibility is $N(u)=N(v)=3$, which is impossible. Thus $3$ is irreducible.

On the other hand,
\[
(1+\sqrt{-5})(1-\sqrt{-5})=6=2\cdot3,
\]
so $3$ divides the product. It divides neither factor: an equality $1\pm\sqrt{-5}=3(a+b\sqrt{-5})$ would require $3a=1$. Hence $3$ is \emph{not} prime. This is the precise failure of unique factorisation that ideals will help us understand later.
\end{solution}

\begin{solution}{ex:zi}
\idea Euclidean division with $N=|\cdot|^2$ works as soon as every point of the plane is at distance $<1$ from the lattice: divide in $\C$, then round.
\writeup (a) Let $z=x+iy$ and let $m,n$ be the nearest integers to $x,y$. Then $|z-(m+in)|^2\le\frac14+\frac14=\frac12$. Given $a$ and $b\ne0$ in $\Z[i]$, take $q$ a Gaussian integer nearest to $a/b$ and $r=a-bq\in\Z[i]$: $N(r)=|b|^2|a/b-q|^2\le\frac12N(b)<N(b)$.

(b) $\Z[\sqrt{-2}]$ is the lattice of points $m+n\sqrt2\,i$. Rounding $x$ to $m$ and $y/\sqrt2$ to $n$ gives $|z-q|^2\le\frac14+2\cdot\frac14=\frac34<1$. The same argument applies.

(c) For $\Z[\sqrt{-5}]$ the same rounding gives only $|z-q|^2\le\frac14+\frac54=\frac32$, which is not $<1$. Indeed the point $z=\frac{1+\sqrt{-5}}2$ is at squared distance $\frac14+\frac54=\frac32$ from the four nearest lattice points $0,1,\sqrt{-5},1+\sqrt{-5}$, and farther from the others. Take $a=1+\sqrt{-5}$, $b=2$, so $a/b=z$: for every $q$, $N(a-bq)=4|z-q|^2\ge6>4=N(b)$. Of course this only shows that this particular $N$ fails; no Euclidean function can exist, since $\Z[\sqrt{-5}]$ is not even a UFD.
\end{solution}

\begin{solution}{ex:zxnotpid}
\idea A generator of $(a,X)$ divides a constant, so is a constant, and divides $X$, so is a unit.
\writeup Let $R$ be a domain that is not a field, $a$ a non-zero non-unit, and suppose $(a,X)=(d)$. Since $d\mid a$ and $a\ne0$ has degree $0$, $d$ is a constant; since $d\mid X$, $X=d\,(e+fX)$ gives $df=1$, so $d$ is a unit and $(a,X)=R[X]$. Then $1=a\,g+X\,h$ for some $g,h$; evaluating at $X=0$, $1=a\,g(0)$, so $a$ is a unit: contradiction. With $R=\Z$, $a=2$: $(2,X)$ is not principal. Conversely, if $R$ is a field, $R[X]$ is Euclidean for the degree, hence a PID. (If $R$ is not a domain, $R[X]$ is not a domain either.)
\end{solution}

\begin{solution}{ex:boolean}
\idea Squaring sums produces the cross terms, and $x^2=x$ forces them to vanish.
\writeup $x+x=(x+x)^2=x^2+2x^2+x^2=4x$, so $2x=0$. In a Boolean domain, $x(x-1)=x^2-x=0$, so $x\in\{0,1\}$: the ring is $\{0,1\}\cong\Z/2\Z$ (and $1\ne0$). If $P$ is a prime ideal, $R/P$ is a Boolean domain, so it is $\Z/2\Z$, a field, and $P$ is maximal. For subsets of $E$ with $A+B=A\triangle B$ and $AB=A\cap B$: $A\cap A=A$, so this ring is Boolean; its ideals of the form $\{A:a\notin A\}$ are maximal, with quotient $\Z/2\Z$ given by ``does $A$ contain $a$?''.
\end{solution}

\begin{solution}{ex:crtinterp}
\idea Find the two ``idempotents'' $e_1\equiv1\bmod X$, $e_1\equiv0\bmod(X-1)^2$ and $e_2$ the other way round; then $P=1\cdot e_1+(\text{target mod }(X-1)^2)\cdot e_2$ reduced mod $X(X-1)^2$.
\writeup The conditions are $P\equiv1\pmod X$ and $P\equiv2(X-1)\pmod{(X-1)^2}$ (a polynomial with $P(1)=0$, $P'(1)=2$ is $2(X-1)$ plus a multiple of $(X-1)^2$, by Taylor's formula). Since $1=(X-1)^2-X(X-2)$, put $e_1=(X-1)^2$ ($\equiv1\bmod X$, $\equiv0\bmod(X-1)^2$) and $e_2=-X(X-2)=1-e_1$. Then $P_0=e_1+2(X-1)e_2$ satisfies both congruences, and the unique solution of degree $\le2$ is its remainder modulo $X(X-1)^2=X^3-2X^2+X$:
\[
\begin{aligned}P_0&=(X-1)^2-2X(X-1)(X-2)\eqbr=-2X^3+7X^2-6X+1\\&\equiv3X^2-4X+1.\end{aligned}
\]
Check: $P=3X^2-4X+1$ has $P(0)=1$, $P(1)=0$, $P'(1)=6-4=2$. (This is Hermite interpolation: conditions on values \emph{and} derivatives are congruences modulo powers of $X-a$.)
\end{solution}

\begin{solution}{ex:x4p1}
\idea Over $\Q$, $X^4+1$ is a cyclotomic polynomial, and Eisenstein applies after a shift. Modulo $p$, three different completions of the square cover all cases, because the product $(-1)\cdot2\cdot(-2)=4$ is a square.
\writeup Over $\Q$: $(X+1)^4+1=X^4+4X^3+6X^2+4X+2$ is Eisenstein at $2$, so it is irreducible (Volume~1, Chapter~7), and so is $X^4+1$ (a factorisation of $X^4+1$ would give one of $(X+1)^4+1$).

Modulo $2$: $X^4+1=(X+1)^4$. Modulo an odd prime $p$: the Legendre symbol is multiplicative (Volume~2), and $\left(\frac{-1}p\right)\left(\frac2p\right)\left(\frac{-2}p\right)=\left(\frac4p\right)=1$, so the three symbols cannot all be $-1$: one of $-1,2,-2$ is a square mod $p$.
\begin{itemize}
\item If $-1\equiv a^2$: $X^4+1=X^4-a^2=(X^2-a)(X^2+a)$.
\item If $2\equiv b^2$: $X^4+1=(X^2+1)^2-2X^2=(X^2+bX+1)(X^2-bX+1)$.
\item If $-2\equiv c^2$: $X^4+1=(X^2-1)^2+2X^2=(X^2-1)^2-(cX)^2=(X^2-cX-1)(X^2+cX-1)$.
\end{itemize}
In each case $X^4+1$ is a product of two quadratics mod $p$. So irreducibility mod $p$ is a sufficient test, never a necessary one: $X^4+1$ is irreducible over $\Q$ but fails the test at every prime. (Chapter~10 explains why: its Galois group $(\Z/2\Z)^2$ has no element of order $4$.)
\end{solution}

\section*{Sources and further reading}

\begin{itemize}
\item M.~Artin, \emph{Algebra}, 2nd ed., Pearson, 2011, Chapters~11 and~12. Rings, ideals and factorisation, with $\Z[\sqrt{-5}]$ as the running example.
\item D.~S.~Dummit and R.~M.~Foote, \emph{Abstract Algebra}, 3rd ed., Wiley, 2004, Chapters~7--9. Including a proof that $\Z[\frac{1+\sqrt{-19}}2]$ is a PID but not Euclidean (Section~8.2).
\item T.~Motzkin, \emph{The Euclidean algorithm}, Bull.\ Amer.\ Math.\ Soc.\ 55 (1949), 1142--1146. The criterion used to show that a ring is not Euclidean for any function.
\item H.~M.~Edwards, \emph{Fermat's Last Theorem: A Genetic Introduction to Algebraic Number Theory}, Springer, 1977. The story of Lamé, Kummer and ideal numbers.
\item J.~von~zur~Gathen and J.~Gerhard, \emph{Modern Computer Algebra}, 3rd ed., Cambridge University Press, 2013. Multi-modular algorithms and polynomial factorisation.
\item L.~Corry, \emph{Steht es alles wirklich schon bei Dedekind? Ideals and factorization between Dedekind and Noether}, in \emph{Ideals and Their Reception}, 2000. The source of Noether's saying.
\end{itemize}
